\documentclass[10pt,a4paper]{amsart}
\usepackage[margin=19mm]{geometry}
\usepackage{amsmath,amssymb,mathtools}
\usepackage[scr=rsfs,cal=euler]{mathalfa}
\usepackage{xfrac}
\usepackage[unicode,hidelinks,bookmarks=false]{hyperref}
\newcommand{\ie}{i.e.\ }
\newcommand{\cf}{cf.\ }
\newcommand{\resp}{respectively\ }
\newcommand{\Subsection}{\subsection}
\providecommand{\nobreakdash}{\nobreak\hbox{-}}
\setlength{\emergencystretch}{2em}
\multlinegap0pt
\usepackage{amssymb}
\usepackage{mathtools}
\usepackage[shortlabels]{enumitem}
\usepackage{xcolor}
\newtheorem{proposition}{Proposition}
\newcommand{\ZZ}{\mathbb Z}
\newcommand{\diag}{\operatorname{diag}}
\newcommand{\tr}{\operatorname{tr}}
\title{Star-Shaped Integral Cartan-Type Matrices and an Egyptian-Fraction Classification of Affine Weighted Trees}
\author{Emilio Torrente-Lujan}
\date{Source: arXiv:2605.23011v1 (CC BY 4.0)}
\begin{document}
\maketitle

\noindent\emph{Source and licence.} Star-Shaped Integral Cartan-Type Matrices and an Egyptian-Fraction Classification of Affine Weighted Trees. Version arXiv:2605.23011v1 (CC BY 4.0), distributed by arXiv under the Creative Commons Attribution 4.0 International licence, http://creativecommons.org/licenses/by/4.0/, as recorded in the arXiv licence metadata for this exact version and shown on the abstract page https://arxiv.org/abs/2605.23011v1. Full licence notice: \texttt{licenses/arxiv-2605.23011-CC-BY-4.0.txt}.\par\smallskip

\noindent\textit{Editorial scope.} Counted unit: the article's proposition eq:determinant (source0.tex lines 151--158). The article's complete original proof of it is delivered with it (source0.tex lines 160--172, one \texttt{proof} environment of 13 lines). No mathematics is written, repaired or re-derived: the statement and the proof are the article's own lines, in the article's own notation, and no retained line is substituted or re-flowed.  Retained uncounted as the source-local context the proof needs: lines 110--146.  Nothing was omitted from any retained span, so the package carries no omission record.  No retained line contains a citation, so the wrapper carries no bibliography and no bibliography entry.  No retained line contains a figure environment, an image-inclusion command or drawing code, so no image is delivered and the package declares no asset dependency.  Editorial formatting only, in addition: an AMS \texttt{amsart} preamble that restates the article's own declarations and macros used by the retained lines, namely the environments \texttt{proposition} on separate counters, together with the standard packages those lines need.  The retained environments print in this document as Proposition 1 in order of appearance, whereas the article prints the same items with the numbering of its own full document.  The complete native source of the article as distributed in the arXiv e-print for this exact version is bundled unchanged as \texttt{sources/201176/source0.tex}.
Let $A_r$ denote the Cartan matrix of type $A_r$,
\[
A_r=\begin{pmatrix}
2&-1&0&\cdots&0\\
-1&2&-1&\ddots&\vdots\\
0&-1&2&\ddots&0\\
\vdots&\ddots&\ddots&\ddots&-1\\
0&\cdots&0&-1&2
\end{pmatrix}.
\]
Fix $m\geq 2$, positive integers $r_1,\ldots,r_m$, and a central weight $k\in\ZZ_{>0}$.  Let $v_i\in\ZZ^{r_i}$ be the column vector with all entries zero except the last entry, which is $-1$.  The star-shaped matrix studied in this paper is
\begin{equation}
B(k;r_1,\ldots,r_m)=
\begin{pmatrix}
A_{r_1} &0&\cdots&0&v_1\\
0&A_{r_2}&\cdots&0&v_2\\
\vdots&\vdots&\ddots&\vdots&\vdots\\
0&0&\cdots&A_{r_m}&v_m\\
v_1^t&v_2^t&\cdots&v_m^t&k
\end{pmatrix}.
\label{eq:star-matrix}
\end{equation}
The underlying graph is a tree with $m$ linear arms attached to one central vertex.  Its dimension is
\begin{equation}
D=1+\sum_{i=1}^m r_i,
\label{eq:dimension}
\end{equation}
and its trace is
\begin{equation}
\tr B=2\sum_{i=1}^m r_i+k.
\end{equation}
The sum of all matrix entries is
\begin{equation}
K^2(B):=\sum_{a,b}B_{ab}=k.
\label{eq:K-square}
\end{equation}
In the previous geometric motivation this number was interpreted as an intersection-theoretic invariant.  In the present matrix/graph treatment it is simply the total entry sum of the weighted star matrix.

\begin{proposition}[Determinant]
For the matrix \eqref{eq:star-matrix},
\begin{equation}
\det B(k;r_1,\ldots,r_m)=
\prod_{i=1}^m (r_i+1)\left(k-m+\sum_{i=1}^m \frac{1}{r_i+1}\right).
\label{eq:determinant}
\end{equation}
\end{proposition}

\begin{proof}
Use the Schur complement with respect to the block diagonal matrix $\diag(A_{r_1},\ldots,A_{r_m})$.  The standard identities
\[
\det A_r=r+1,\qquad (A_r^{-1})_{rr}=\frac{r}{r+1}
\]
give
\begin{align*}
\det B
&=\prod_{i=1}^m \det A_{r_i}\left(k-\sum_{i=1}^m v_i^t A_{r_i}^{-1}v_i\right)\\
&=\prod_{i=1}^m (r_i+1)\left(k-\sum_{i=1}^m\frac{r_i}{r_i+1}\right)\\
&=\prod_{i=1}^m (r_i+1)\left(k-m+\sum_{i=1}^m\frac{1}{r_i+1}\right).
\end{align*}
\end{proof}

\end{document}
